\documentclass[10pt,a4paper]{amsart}
\usepackage[margin=19mm]{geometry}
\usepackage{amsmath,amssymb,mathtools}
\usepackage[scr=rsfs,cal=euler]{mathalfa}
\usepackage{xfrac}
\usepackage[unicode,hidelinks,bookmarks=false]{hyperref}
\newcommand{\ie}{i.e.\ }
\newcommand{\cf}{cf.\ }
\newcommand{\resp}{respectively\ }
\newcommand{\Subsection}{\subsection}
\providecommand{\nobreakdash}{\nobreak\hbox{-}}
\setlength{\emergencystretch}{2em}
\multlinegap0pt
\usepackage{amssymb}
\usepackage{mathtools}
\newtheorem{lemma}{Lemma}
\newcommand{\bin}[2]{\binom{#1}{#2}}
\title{The minimum of the graph likelihood}
\author{Simone Severini, Eric W. Weisstein}
\date{Source: arXiv:2608.19467v1 (CC BY 4.0)}
\begin{document}
\maketitle

\noindent\emph{Source and licence.} The minimum of the graph likelihood. Version arXiv:2608.19467v1 (CC BY 4.0), distributed by arXiv under the Creative Commons Attribution 4.0 International licence, http://creativecommons.org/licenses/by/4.0/, as recorded in the arXiv licence metadata for this exact version and shown on the abstract page https://arxiv.org/abs/2608.19467v1. Full licence notice: \texttt{licenses/arxiv-2608.19467-CC-BY-4.0.txt}.\par\smallskip

\noindent\textit{Editorial scope.} Counted unit: the article's lemma lem:greedy (source0.tex lines 657--659). The article's complete original proof of it is delivered with it (source0.tex lines 661--681, one \texttt{proof} environment of 21 lines). No mathematics is written, repaired or re-derived: the statement and the proof are the article's own lines, in the article's own notation, and no retained line is substituted or re-flowed.  No further source-local context is needed: every cross-reference of every retained line resolves inside the two delivered spans.  Nothing was omitted from any retained span, so the package carries no omission record.  No retained line contains a citation, so the wrapper carries no bibliography and no bibliography entry.  No retained line contains a figure environment, an image-inclusion command or drawing code, so no image is delivered and the package declares no asset dependency.  Editorial formatting only, in addition: an AMS \texttt{amsart} preamble that restates the article's own declarations and macros used by the retained lines, namely the environments \texttt{lemma} on separate counters, together with the standard packages those lines need.  The retained environments print in this document as Lemma 1 in order of appearance, whereas the article prints the same items with the numbering of its own full document.  The complete native source of the article as distributed in the arXiv e-print for this exact version is bundled unchanged as \texttt{sources/208092/source0.tex}.
\begin{lemma}\label{lem:greedy}
$W(m)=\sum_{j=0}^{m-1}\log_2\bin{m+j}{j}$, attained by the path that empties one side first.
\end{lemma}

\begin{proof}
Let $f(a,b)$ denote the minimum cost of a deletion path from $(a,b)$, with $f(a,0)=f(0,b)=0$. We show by induction on $a+b$ that for $a\le b$,
\[
f(a,b)=S(a,b):=\sum_{j=0}^{a-1}\log_2\bin{b+j}{j},
\]
which is the cost of emptying the first side; by symmetry this suffices. The cases $a=0$ and $a+b\le2$ are clear. Let $1\le a\le b$. Emptying the first side has cost $\log_2\bin{a+b-1}{b}+S(a-1,b)=S(a,b)$, since $\bin{a+b-1}{b}=\bin{a+b-1}{a-1}$. So $f(a,b)\le S(a,b)$, and it remains to exclude the other first step, that is to show
\begin{equation}\label{eq:exchange}
S(a,b)\;\le\; \log_2\bin{a+b-1}{a}+f(a,b-1).
\end{equation}
Suppose first $a\le b-1$, so that $f(a,b-1)=S(a,b-1)$ by induction. Since
$\bin{b+j}{j}\big/\bin{b-1+j}{j}=(b+j)/b$, we get
\begin{align*}
S(a,b)-S(a,b-1)
&=\sum_{j=0}^{a-1}\log_2\frac{b+j}{b}\\
&=\log_2\frac{\prod_{j=0}^{a-1}(b+j)}{b^{a}},\\
\log_2\bin{a+b-1}{a}
&=\log_2\frac{\prod_{j=0}^{a-1}(b+j)}{a!},
\end{align*}
so \eqref{eq:exchange} is equivalent to $a!\le b^{a}$, which holds because $a\le b$. If instead $a=b$, then $f(a,b-1)=S(a-1,a)$ by induction and
$S(a,a)-S(a-1,a)=\log_2\bin{2a-1}{a-1}=\log_2\bin{2a-1}{a}$, so \eqref{eq:exchange} holds with equality, as symmetry demands.
\end{proof}

\end{document}
